\begin{abstract}
This paper studies daily closing-price forecasting for major Indonesian banking stocks using recurrent sequence models, macroeconomic regressors, and statistical validation. The primary ablation compares LSTM and BiLSTM under two feature settings (\texttt{Without\_Macro} and \texttt{With\_Macro}) and two optimization modes (default vs.\ budgeted tuning via Genetic Algorithm), resulting in eight experimental cases per bank. To improve rigor, we run repeated multi-seed evaluations across seven banks and apply Wilcoxon signed-rank, Diebold-Mariano, and SHAP-based analyses. Empirically, BiLSTM-based configurations dominate most banks, with best-case RMSE improvements over the baseline ranging from 4.91\% to 35.66\% for six out of seven banks. As an extension, we further evaluate Informer/Autoformer Transformer models and recurrent variants based on normalization, convolution, and temporal attention. The Transformer extension shows strong aggregate performance, with \texttt{informer\_data} achieving RMSE 401.02, MAE 325.11, and MAPE 5.94\%, while recurrent variants do not outperform the original recurrent benchmark on average. These extension results provide promising directions for future statistically matched validation.
\end{abstract}

\renewcommand\IEEEkeywordsname{Keywords}
\begin{IEEEkeywords}
stock forecasting, LSTM, BiLSTM, Informer, Autoformer, macroeconomic features, genetic algorithm, Wilcoxon test, DM test, SHAP
\end{IEEEkeywords}
